\textbf{Proof.}
Write \(L=\log(en)=1+\log n\).  We first prove the lower-dimensional band directly and then invoke the cited known-\(\mathbf Q\) theorem in the complementary band.

\emph{Paired-sign band.}
Let \(k=2\lfloor d/2\rfloor\), so that \(k\) is even and \(d/2\le k\le d\).  Define \(\mathbf Q^{(0)}\in\Delta_d\) by \(Q^{(0)}_x=1/k\) for \(x\le k\) and \(Q^{(0)}_x=0\) otherwise.  For \(\sigma=(\sigma_1,\ldots,\sigma_{k/2})\in\{-1,1\}^{k/2}\) and \(0<\alpha\le1/2\), define \(\mathbf P_\sigma\in\Delta_d\) by
\[
 P_{\sigma,2j-1}=\frac{1+\alpha\sigma_j}{k},\qquad
 P_{\sigma,2j}=\frac{1-\alpha\sigma_j}{k},
 \qquad 1\le j\le k/2,
\]
and set its remaining coordinates to zero.  Nonnegativity follows from \(\alpha\le1/2\), and summing within pairs shows that the coordinates sum to one.  Furthermore,
\[
 \lVert\mathbf P_\sigma-\mathbf Q^{(0)}\rVert_1
 =\sum_{j=1}^{k/2}\left(\frac{\alpha}{k}+\frac{\alpha}{k}\right)
 =\alpha .
\]

Let \(M\) be the uniform mixture of the laws \(\mathbf P_\sigma^{\otimes n}\), and let \(Q_0=(\mathbf Q^{(0)})^{\otimes n}\).  If \(Z\sim\mathbf Q^{(0)}\), direct multiplication of the one-observation likelihood ratios gives, for every \(\sigma,\tau\),
\[
 E\left[
 \frac{P_\sigma(Z)}{Q^{(0)}(Z)}
 \frac{P_\tau(Z)}{Q^{(0)}(Z)}
 \right]
 =1+\frac{2\alpha^2}{k}\sum_{j=1}^{k/2}\sigma_j\tau_j.
\]
Independence of the \(n\) observations and expansion of the squared mixture likelihood ratio therefore yield
\[
 1+\chi^2(M,Q_0)
 =E_{\sigma,\tau}\left(
 1+\frac{2\alpha^2}{k}\sum_{j=1}^{k/2}\sigma_j\tau_j
 \right)^n.
\]
Using \(1+t\le e^t\), independence of the Rademacher products \(\sigma_j\tau_j\), and \(\cosh t\le e^{t^2/2}\), we obtain
\[
 1+\chi^2(M,Q_0)
 \le
 \left\{\cosh\left(\frac{2n\alpha^2}{k}\right)\right\}^{k/2}
 \le \exp\left(\frac{n^2\alpha^4}{k}\right).
\]
Choose a universal \(b>0\) sufficiently small that
\[
 \frac12\sqrt{e^{b^2}-1}\le\frac12,
\]
and set \(\alpha^2=b\sqrt{k}/n\).  Whenever \(d\le K L^2\), for any fixed universal \(K<\infty\), this choice satisfies \(\alpha\le1/2\) once \(n\) exceeds a universal threshold depending only on \(K\).  The preceding display then gives \(\chi^2(M,Q_0)\le e^{b^2}-1\).  Cauchy--Schwarz applied to the mixture likelihood ratio gives
\[
 \lVert M-Q_0\rVert_{\mathrm{TV}}
 \le\frac12\sqrt{\chi^2(M,Q_0)}\le\frac12.
\]

We spell out the risk reduction.  For an arbitrary estimator \(\widehat L\), let \(\phi=\mathbf 1\{\widehat L\ge\alpha/2\}\).  Under the null pair \((\mathbf Q^{(0)},\mathbf Q^{(0)})\), the target is zero, whereas under every alternative pair \((\mathbf P_\sigma,\mathbf Q^{(0)})\), it is \(\alpha\).  The second sample has the same law under the null and all alternatives, so tensoring both \(M\) and \(Q_0\) with that common law leaves their total-variation distance unchanged.  Consequently,
\[
 \begin{aligned}
 &\max\left\{
 E_{\mathbf Q^{(0)},\mathbf Q^{(0)}}\widehat L^2,
 \sup_\sigma E_{\mathbf P_\sigma,\mathbf Q^{(0)}}(\widehat L-\alpha)^2
 \right\}\\
 &\quad\ge
 \frac{\alpha^2}{8}
 \left\{Q_0(\phi=1)+M(\phi=0)\right\}
 \ge \frac{\alpha^2}{8}
 \left\{1-\lVert M-Q_0\rVert_{\mathrm{TV}}\right\}
 \ge\frac{\alpha^2}{16}.
 \end{aligned}
\]
Taking the infimum over \(\widehat L\) proves
\[
 \mathfrak R^{L_1}_{n,d}\ge c\frac{\sqrt{k}}{n}.
\]
If \(d\le K L^2\), then \(k\ge d/2\) and hence
\[
 \frac{\sqrt{k}}n
 \ge \frac{d}{\sqrt{2K}\,nL}.
\]
Thus the asserted lower bound holds throughout this band.

\emph{Uniform known-\(\mathbf Q\) band.}
Apply \(\mathrm{lem:jiao-han-weissman-known-q-lower}\) with its logarithmic constant fixed at \(1/2\), and denote the corresponding constants by \(C'>0\) and \(c'>0\).  Take \(\mathbf Q^{(u)}=(1/d,\ldots,1/d)\).  We will choose \(\rho>0\) small enough that, for all sufficiently large \(n\),
\[
 d\le\rho nL\quad\Longrightarrow\quad d\le n\log n.
\]
For such \(n,d\), the minimum in the cited theorem is attained by its square-root term, and therefore
\[
 T_{n,d}(\mathbf Q^{(u)})
 =d\sqrt{\frac{1/d}{n\log n}}
 =\sqrt{\frac d{n\log n}}.
\]
Choose the universal constant \(K\) large enough that \(d\ge K L^2\) implies
\[
 T_{n,d}(\mathbf Q^{(u)})
 \ge C'\sqrt{\frac{\log n}{n}}.
\]
Moreover, identically in \(d\),
\[
 \frac{T_{n,d}(\mathbf Q^{(u)})}{\sqrt d\log n/n}
 =\frac{\sqrt n}{(\log n)^{3/2}},
\]
which is at least \(C'\) for all sufficiently large \(n\).  Finally, if \(d\le\rho nL\), then \(\log d\le2\log n\) for all sufficiently large \(n\), so the cited requirement \(\log n\ge\frac12\log d\) also holds.  All hypotheses of the cited lower bound are therefore satisfied when \(d\ge K L^2\), and its fixed-multinomial comparison gives
\[
 \inf_{\widehat L}
 \sup_{\mathbf P\in\Delta_d}
 E_{\mathbf P,\mathbf Q^{(u)}}
 \left(\widehat L-\lVert\mathbf P-\mathbf Q^{(u)}\rVert_1\right)^2
 \ge c'\frac d{n\log n}
 \ge c'\frac d{nL}.
\]
The unrestricted two-sample minimax risk cannot be smaller: restricting its supremum to \(\mathbf Q=\mathbf Q^{(u)}\) and then revealing this fixed \(\mathbf Q^{(u)}\) only makes the estimation problem easier.  Hence the same lower bound holds for \(\mathfrak R^{L_1}_{n,d}\).

It remains only to choose constants uniformly.  First choose \(K\) as above, and then choose a universal integer \(n_0\) large enough for every displayed large-\(n\) condition, including \(b\sqrt{k}/n\le1/4\) when \(k\le K L^2\).  Choose \(\rho>0\) so small that the uniform-band implications hold for \(n\ge n_0\) and, in addition,
\[
 \rho n\log(en)<2\qquad(1\le n<n_0).
\]
For \(n<n_0\) the claimed range \(2\le d\le\rho n\log(en)\) is then empty.  For \(n\ge n_0\), either \(d\le K L^2\) and the paired-sign argument applies, or \(d\ge K L^2\) and the cited uniform argument applies.  Decreasing \(c>0\) to the smaller of the two universal band constants and taking \(C=K\) proves the result. \(\square\)